\chapter*{Abstract}
\noindent
Corporate bankruptcy prediction is complicated by severe class imbalance, and
much of the recent literature addresses this by synthesising minority
observations, reporting accuracies near $99\%$ on data that is partly fabricated.
This thesis takes the opposite approach: it trains exclusively on real
historical firms, restoring balance through cluster-centroid undersampling and
evaluating models by the recall-oriented $F_2$-score. On the Taiwanese
Bankruptcy Dataset ($6{,}819$ firms, $95$ features, $3.23\%$ bankrupt), three
models of increasing complexity are compared. Gradient boosting attains the best
cross-validated performance ($F_2 = 0.550$), ahead of a random forest ($0.546$)
and a logistic-regression baseline ($0.482$), and competitive with the
undersampling benchmark of Wang and Liu (2021) while introducing no synthetic
data. The central contribution is an interpretability layer that their work
lacked: comparing native, permutation, and SHAP importance across all three
models reveals a systematic divergence. Native impurity importance is distorted
by a known bias towards continuous features, permutation importance collapses
under the collinearity of the ratios, and SHAP --- the only method yielding
signed, comparable attributions for every model --- provides the most
trustworthy ranking, dominated by leverage and debt-servicing measures. The
thesis argues that this comparison, not any single importance list, is what
widens the black box.

\vspace{1em}
\noindent\textit{\textbf{Keywords: }bankruptcy prediction, class imbalance,
undersampling, feature importance, SHAP}
